% --- Archon physics formalization source begin ---
% archon:physics
% archon:covers PhyXMiniProblems/problem_phyx_mini_0508.lean
% archon:source-report reports/phyx_mini/problem_phyx_mini_0508.source.json

\chapter{Physics problem phyx\_mini\_0508}
\label{ch:PhyXMiniProblems_problem_phyx_mini_0508}

\paragraph{Problem source.}
\#\# Physical scenario

Figure shows the wave function of an electron in a rigid box. The electron energy is \textbackslash{}( 6.0 \textbackslash{}text\{ eV\} \textbackslash{}).

\#\# Question

How long is the box?

\#\# Answer choices

- A: A: \textbackslash{}( 0.5 \textbackslash{}text\{ nm\} \textbackslash{})

- B: B: \textbackslash{}( 3.0 \textbackslash{}text\{ nm\} \textbackslash{})

- C: C: \textbackslash{}( 1.0 \textbackslash{}text\{ nm\} \textbackslash{})

- D: D: \textbackslash{}( 2.0 \textbackslash{}text\{ nm\} \textbackslash{})

\#\# Figure caption (auxiliary; use the image as primary evidence)

The image is a graph depicting a wave function. It features a sinusoidal wave illustrated in red, cutting across and oscillating around the x-axis, which is horizontal. The wave starts slightly above the x-axis, has two full oscillations (peaks and troughs), and ends slightly above the x-axis again.

On the left side, the y-axis is vertical and labeled with the function \textbackslash{}(\textbackslash{}psi(x)\textbackslash{}). The x-axis is labeled with \textbackslash{}(x\textbackslash{}). The wave function is shown as symmetric in nature, extending across the x-axis. There are no numerical values marked on the axes.

\#\# Dataset metadata

Quantum Phenomena; Physical Model Grounding Reasoning, Spatial Relation Reasoning

\paragraph{Recorded answer/context.}
C: C: \textbackslash{}( 1.0 \textbackslash{}text\{ nm\} \textbackslash{})

\paragraph{Figure/image path.}
/root/proposal\_for\_physic/science-mango/phyx\_mini\_run/phyx\_data/test\_image/508.png

\paragraph{Formalization target.}
create a compiling Lean file with sorry bodies at `PhyXMiniProblems/problem\_phyx\_mini\_0508.lean`. The Lean declarations must preserve the physical quantities, dimensions or dimensional roles, figure labels, governing-law hypotheses, and final relation expressed by this problem.
Use Mathlib/Physlib names found through LeanExplore where available. If a physics API is missing, introduce faithful local abstractions rather than scalar placeholder aliases.

\begin{theorem}[Physics formalization target]
\label{thm:physics:phyx_mini_0508:target}
\lean{PhyXMiniProblems.ProblemPhyXMini0508.problem_phyx_mini_0508}
\uses{def:physics:phyx-mini-0508:phyxminiproblems-problemphyxmini0508-lengthinnanometers, def:physics:phyx-mini-0508:phyxminiproblems-problemphyxmini0508-rigidboxelectronsetup, def:physics:phyx-mini-0508:phyxminiproblems-problemphyxmini0508-matchesrigidboxelectronscenario, def:physics:phyx-mini-0508:phyxminiproblems-problemphyxmini0508-matchesproblemenergyreadout, def:physics:phyx-mini-0508:phyxminiproblems-problemphyxmini0508-useselectronmassreferencedata, def:physics:phyx-mini-0508:phyxminiproblems-problemphyxmini0508-matchessuppliedwavefunctionfigure, def:physics:phyx-mini-0508:phyxminiproblems-problemphyxmini0508-representsrigidboxstationarymode, def:physics:phyx-mini-0508:phyxminiproblems-problemphyxmini0508-satisfiesinfiniterigidboxenergylaw, def:physics:phyx-mini-0508:phyxminiproblems-problemphyxmini0508-hasphysicalrigidboxparameters, def:physics:phyx-mini-0508:phyxminiproblems-problemphyxmini0508-recordeddatasetanswer, def:physics:phyx-mini-0508:phyxminiproblems-problemphyxmini0508-roundedtonearesttenth, def:physics:phyx-mini-0508:phyxminiproblems-problemphyxmini0508-isuniquematchingdisplayedlength}
The assigned autoformalize agent should translate this physics problem into Lean declarations in the covered file, with theorem and lemma proof bodies written as `by sorry`.
\end{theorem}
\begin{proof}
This is an autoformalization task, not a proof task. Produce statements that can later be proved by the physics prover without weakening the physical model.
\end{proof}

% --- Archon declaration topology begin ---
\section{Declaration topology}
\begin{definition}[Length Quantity]
  \label{def:physics:phyx-mini-0508:phyxminiproblems-problemphyxmini0508-lengthquantity}
  \lean{PhyXMiniProblems.ProblemPhyXMini0508.LengthQuantity}
  A nonnegative, unit-independent physical length
\end{definition}
\begin{proof}
  This is the declared model definition of Length Quantity.
\end{proof}
\begin{definition}[Mass Quantity]
  \label{def:physics:phyx-mini-0508:phyxminiproblems-problemphyxmini0508-massquantity}
  \lean{PhyXMiniProblems.ProblemPhyXMini0508.MassQuantity}
  A nonnegative, unit-independent physical mass
\end{definition}
\begin{proof}
  This is the declared model definition of Mass Quantity.
\end{proof}
\begin{definition}[length Readout]
  \label{def:physics:phyx-mini-0508:phyxminiproblems-problemphyxmini0508-lengthreadout}
  \lean{PhyXMiniProblems.ProblemPhyXMini0508.lengthReadout}
  \uses{def:physics:phyx-mini-0508:phyxminiproblems-problemphyxmini0508-lengthquantity}
  Read a physical length in a selected length unit
\end{definition}
\begin{proof}
  This is the declared model definition of length Readout.
\end{proof}
\begin{definition}[length In Meters]
  \label{def:physics:phyx-mini-0508:phyxminiproblems-problemphyxmini0508-lengthinmeters}
  \lean{PhyXMiniProblems.ProblemPhyXMini0508.lengthInMeters}
  \uses{def:physics:phyx-mini-0508:phyxminiproblems-problemphyxmini0508-lengthquantity, def:physics:phyx-mini-0508:phyxminiproblems-problemphyxmini0508-lengthreadout}
  SI metre readout of a physical length
\end{definition}
\begin{proof}
  This is the declared model definition of length In Meters.
\end{proof}
\begin{definition}[length In Nanometers]
  \label{def:physics:phyx-mini-0508:phyxminiproblems-problemphyxmini0508-lengthinnanometers}
  \lean{PhyXMiniProblems.ProblemPhyXMini0508.lengthInNanometers}
  \uses{def:physics:phyx-mini-0508:phyxminiproblems-problemphyxmini0508-lengthquantity, def:physics:phyx-mini-0508:phyxminiproblems-problemphyxmini0508-lengthreadout}
  Nanometre readout used by the displayed answers
\end{definition}
\begin{proof}
  This is the declared model definition of length In Nanometers.
\end{proof}
\begin{definition}[mass In Kilograms]
  \label{def:physics:phyx-mini-0508:phyxminiproblems-problemphyxmini0508-massinkilograms}
  \lean{PhyXMiniProblems.ProblemPhyXMini0508.massInKilograms}
  \uses{def:physics:phyx-mini-0508:phyxminiproblems-problemphyxmini0508-massquantity}
  SI kilogram readout of a physical mass
\end{definition}
\begin{proof}
  This is the declared model definition of mass In Kilograms.
\end{proof}
\begin{definition}[energy In Joules]
  \label{def:physics:phyx-mini-0508:phyxminiproblems-problemphyxmini0508-energyinjoules}
  \lean{PhyXMiniProblems.ProblemPhyXMini0508.energyInJoules}
  SI joule readout of a physical energy
\end{definition}
\begin{proof}
  This is the declared model definition of energy In Joules.
\end{proof}
\begin{definition}[energy In Electron Volts]
  \label{def:physics:phyx-mini-0508:phyxminiproblems-problemphyxmini0508-energyinelectronvolts}
  \lean{PhyXMiniProblems.ProblemPhyXMini0508.energyInElectronVolts}
  \uses{def:physics:phyx-mini-0508:phyxminiproblems-problemphyxmini0508-energyinjoules}
  Electron-volt readout, grounded in Physlib's calibrated electron volt
\end{definition}
\begin{proof}
  This is the declared model definition of energy In Electron Volts.
\end{proof}
\begin{definition}[Particle Species]
  \label{def:physics:phyx-mini-0508:phyxminiproblems-problemphyxmini0508-particlespecies}
  \lean{PhyXMiniProblems.ProblemPhyXMini0508.ParticleSpecies}
  Particle species represented by the state
\end{definition}
\begin{proof}
  This is the declared model definition of Particle Species.
\end{proof}
\begin{definition}[Confinement Model]
  \label{def:physics:phyx-mini-0508:phyxminiproblems-problemphyxmini0508-confinementmodel}
  \lean{PhyXMiniProblems.ProblemPhyXMini0508.ConfinementModel}
  Confining potential used by the model
\end{definition}
\begin{proof}
  This is the declared model definition of Confinement Model.
\end{proof}
\begin{definition}[Figure Label]
  \label{def:physics:phyx-mini-0508:phyxminiproblems-problemphyxmini0508-figurelabel}
  \lean{PhyXMiniProblems.ProblemPhyXMini0508.FigureLabel}
  Labels printed on the two axes in the supplied raster
\end{definition}
\begin{proof}
  This is the declared model definition of Figure Label.
\end{proof}
\begin{definition}[Lobe Sign]
  \label{def:physics:phyx-mini-0508:phyxminiproblems-problemphyxmini0508-lobesign}
  \lean{PhyXMiniProblems.ProblemPhyXMini0508.LobeSign}
  The sign of one consecutive lobe of the plotted real wave amplitude
\end{definition}
\begin{proof}
  This is the declared model definition of Lobe Sign.
\end{proof}
\begin{definition}[Rigid Box Wave Function Figure]
  \label{def:physics:phyx-mini-0508:phyxminiproblems-problemphyxmini0508-rigidboxwavefunctionfigure}
  \lean{PhyXMiniProblems.ProblemPhyXMini0508.RigidBoxWaveFunctionFigure}
  \uses{def:physics:phyx-mini-0508:phyxminiproblems-problemphyxmini0508-figurelabel, def:physics:phyx-mini-0508:phyxminiproblems-problemphyxmini0508-lobesign}
  Qualitative and integer-valued evidence read from the primary image. The flat red pieces before and after the oscillatory segment lie on the horizontal axis; the raster does not mark numerical coordinate values or draw explicit vertical box walls
\end{definition}
\begin{proof}
  This is the declared model definition of Rigid Box Wave Function Figure.
\end{proof}
\begin{definition}[Rigid Box Electron Setup]
  \label{def:physics:phyx-mini-0508:phyxminiproblems-problemphyxmini0508-rigidboxelectronsetup}
  \lean{PhyXMiniProblems.ProblemPhyXMini0508.RigidBoxElectronSetup}
  \uses{def:physics:phyx-mini-0508:phyxminiproblems-problemphyxmini0508-lengthquantity, def:physics:phyx-mini-0508:phyxminiproblems-problemphyxmini0508-massquantity, def:physics:phyx-mini-0508:phyxminiproblems-problemphyxmini0508-particlespecies, def:physics:phyx-mini-0508:phyxminiproblems-problemphyxmini0508-confinementmodel, def:physics:phyx-mini-0508:phyxminiproblems-problemphyxmini0508-rigidboxwavefunctionfigure}
  Independent data for the electron state and the rigid box. The box length is an observable field, not a definition involving the recorded answer. waveRepresentative is retained because an L² state is an almost-everywhere equivalence class whereas the plotted sine profile is pointwise
\end{definition}
\begin{proof}
  This is the declared model definition of Rigid Box Electron Setup.
\end{proof}
\begin{definition}[Matches Rigid Box Electron Scenario]
  \label{def:physics:phyx-mini-0508:phyxminiproblems-problemphyxmini0508-matchesrigidboxelectronscenario}
  \lean{PhyXMiniProblems.ProblemPhyXMini0508.MatchesRigidBoxElectronScenario}
  \uses{def:physics:phyx-mini-0508:phyxminiproblems-problemphyxmini0508-rigidboxelectronsetup}
  Categorical facts explicitly stated in the prose
\end{definition}
\begin{proof}
  This is the declared model definition of Matches Rigid Box Electron Scenario.
\end{proof}
\begin{definition}[Matches Problem Energy Readout]
  \label{def:physics:phyx-mini-0508:phyxminiproblems-problemphyxmini0508-matchesproblemenergyreadout}
  \lean{PhyXMiniProblems.ProblemPhyXMini0508.MatchesProblemEnergyReadout}
  \uses{def:physics:phyx-mini-0508:phyxminiproblems-problemphyxmini0508-energyinelectronvolts, def:physics:phyx-mini-0508:phyxminiproblems-problemphyxmini0508-rigidboxelectronsetup}
  The energy value explicitly supplied by the problem text
\end{definition}
\begin{proof}
  This is the declared model definition of Matches Problem Energy Readout.
\end{proof}
\begin{definition}[Uses Electron Mass Reference Data]
  \label{def:physics:phyx-mini-0508:phyxminiproblems-problemphyxmini0508-useselectronmassreferencedata}
  \lean{PhyXMiniProblems.ProblemPhyXMini0508.UsesElectronMassReferenceData}
  \uses{def:physics:phyx-mini-0508:phyxminiproblems-problemphyxmini0508-massinkilograms, def:physics:phyx-mini-0508:phyxminiproblems-problemphyxmini0508-rigidboxelectronsetup}
  Reference mass data needed to evaluate the infinite-well law. Physlib has no electron-mass constant in the searched API, so the 2022 CODATA kilogram readout is stated separately from the problem and figure readouts
\end{definition}
\begin{proof}
  This is the declared model definition of Uses Electron Mass Reference Data.
\end{proof}
\begin{definition}[Matches Supplied Wave Function Figure]
  \label{def:physics:phyx-mini-0508:phyxminiproblems-problemphyxmini0508-matchessuppliedwavefunctionfigure}
  \lean{PhyXMiniProblems.ProblemPhyXMini0508.MatchesSuppliedWaveFunctionFigure}
  \uses{def:physics:phyx-mini-0508:phyxminiproblems-problemphyxmini0508-rigidboxwavefunctionfigure}
  Raw qualitative facts visible in the supplied wavefunction raster
\end{definition}
\begin{proof}
  This is the declared model definition of Matches Supplied Wave Function Figure.
\end{proof}
\begin{definition}[Represents Rigid Box Stationary Mode]
  \label{def:physics:phyx-mini-0508:phyxminiproblems-problemphyxmini0508-representsrigidboxstationarymode}
  \lean{PhyXMiniProblems.ProblemPhyXMini0508.RepresentsRigidBoxStationaryMode}
  \uses{def:physics:phyx-mini-0508:phyxminiproblems-problemphyxmini0508-lengthinmeters, def:physics:phyx-mini-0508:phyxminiproblems-problemphyxmini0508-rigidboxelectronsetup}
  The plotted profile is interpreted as the nth stationary mode of the rigid box. In an infinite well, the mode number equals the number of half-wave lobes. The analytic representative vanishes outside [0,L] and has sine profile A sin(n π x / L) inside. This relation contains no numerical value for L
\end{definition}
\begin{proof}
  This is the declared model definition of Represents Rigid Box Stationary Mode.
\end{proof}
\begin{definition}[Satisfies Infinite Rigid Box Energy Law]
  \label{def:physics:phyx-mini-0508:phyxminiproblems-problemphyxmini0508-satisfiesinfiniterigidboxenergylaw}
  \lean{PhyXMiniProblems.ProblemPhyXMini0508.SatisfiesInfiniteRigidBoxEnergyLaw}
  \uses{def:physics:phyx-mini-0508:phyxminiproblems-problemphyxmini0508-lengthinmeters, def:physics:phyx-mini-0508:phyxminiproblems-problemphyxmini0508-massinkilograms, def:physics:phyx-mini-0508:phyxminiproblems-problemphyxmini0508-energyinjoules, def:physics:phyx-mini-0508:phyxminiproblems-problemphyxmini0508-rigidboxelectronsetup}
  The one-dimensional infinite-square-well spectrum Eₙ = (n π ℏ)² / (2 m L²). The equation is stated in coherent SI readouts. Constants.ℏ is Physlib's reduced Planck constant in joule-seconds. This is a general governing law and does not assume the requested length or an answer choice
\end{definition}
\begin{proof}
  This is the declared model definition of Satisfies Infinite Rigid Box Energy Law.
\end{proof}
\begin{definition}[Has Physical Rigid Box Parameters]
  \label{def:physics:phyx-mini-0508:phyxminiproblems-problemphyxmini0508-hasphysicalrigidboxparameters}
  \lean{PhyXMiniProblems.ProblemPhyXMini0508.HasPhysicalRigidBoxParameters}
  \uses{def:physics:phyx-mini-0508:phyxminiproblems-problemphyxmini0508-lengthinmeters, def:physics:phyx-mini-0508:phyxminiproblems-problemphyxmini0508-massinkilograms, def:physics:phyx-mini-0508:phyxminiproblems-problemphyxmini0508-energyinjoules, def:physics:phyx-mini-0508:phyxminiproblems-problemphyxmini0508-rigidboxelectronsetup}
  Positivity conditions selecting a nondegenerate physical setup
\end{definition}
\begin{proof}
  This is the declared model definition of Has Physical Rigid Box Parameters.
\end{proof}
\begin{definition}[Answer Choice]
  \label{def:physics:phyx-mini-0508:phyxminiproblems-problemphyxmini0508-answerchoice}
  \lean{PhyXMiniProblems.ProblemPhyXMini0508.AnswerChoice}
  The four answer labels printed with the problem
\end{definition}
\begin{proof}
  This is the declared model definition of Answer Choice.
\end{proof}
\begin{definition}[displayed Length In Nanometers]
  \label{def:physics:phyx-mini-0508:phyxminiproblems-problemphyxmini0508-displayedlengthinnanometers}
  \lean{PhyXMiniProblems.ProblemPhyXMini0508.displayedLengthInNanometers}
  \uses{def:physics:phyx-mini-0508:phyxminiproblems-problemphyxmini0508-answerchoice}
  Nanometre value printed beside each answer label
\end{definition}
\begin{proof}
  This is the declared model definition of displayed Length In Nanometers.
\end{proof}
\begin{definition}[recorded Dataset Answer]
  \label{def:physics:phyx-mini-0508:phyxminiproblems-problemphyxmini0508-recordeddatasetanswer}
  \lean{PhyXMiniProblems.ProblemPhyXMini0508.recordedDatasetAnswer}
  \uses{def:physics:phyx-mini-0508:phyxminiproblems-problemphyxmini0508-answerchoice}
  The answer label recorded by the supplied dataset
\end{definition}
\begin{proof}
  This is the declared model definition of recorded Dataset Answer.
\end{proof}
\begin{definition}[rounded To Nearest Tenth]
  \label{def:physics:phyx-mini-0508:phyxminiproblems-problemphyxmini0508-roundedtonearesttenth}
  \lean{PhyXMiniProblems.ProblemPhyXMini0508.roundedToNearestTenth}
  Round a nanometre readout to the tenth-nanometre precision shown
\end{definition}
\begin{proof}
  This is the declared model definition of rounded To Nearest Tenth.
\end{proof}
\begin{definition}[Matches Displayed Length]
  \label{def:physics:phyx-mini-0508:phyxminiproblems-problemphyxmini0508-matchesdisplayedlength}
  \lean{PhyXMiniProblems.ProblemPhyXMini0508.MatchesDisplayedLength}
  \uses{def:physics:phyx-mini-0508:phyxminiproblems-problemphyxmini0508-lengthinnanometers, def:physics:phyx-mini-0508:phyxminiproblems-problemphyxmini0508-rigidboxelectronsetup, def:physics:phyx-mini-0508:phyxminiproblems-problemphyxmini0508-answerchoice, def:physics:phyx-mini-0508:phyxminiproblems-problemphyxmini0508-displayedlengthinnanometers, def:physics:phyx-mini-0508:phyxminiproblems-problemphyxmini0508-roundedtonearesttenth}
  A choice agrees with the calculated length at the displayed precision
\end{definition}
\begin{proof}
  This is the declared model definition of Matches Displayed Length.
\end{proof}
\begin{definition}[Is Unique Matching Displayed Length]
  \label{def:physics:phyx-mini-0508:phyxminiproblems-problemphyxmini0508-isuniquematchingdisplayedlength}
  \lean{PhyXMiniProblems.ProblemPhyXMini0508.IsUniqueMatchingDisplayedLength}
  \uses{def:physics:phyx-mini-0508:phyxminiproblems-problemphyxmini0508-rigidboxelectronsetup, def:physics:phyx-mini-0508:phyxminiproblems-problemphyxmini0508-answerchoice, def:physics:phyx-mini-0508:phyxminiproblems-problemphyxmini0508-matchesdisplayedlength}
  The selected label is the unique displayed choice matching the length
\end{definition}
\begin{proof}
  This is the declared model definition of Is Unique Matching Displayed Length.
\end{proof}
\begin{lemma}[mode Number eq four]
  \label{lem:physics:phyx-mini-0508:phyxminiproblems-problemphyxmini0508-modenumber-eq-four}
  \lean{PhyXMiniProblems.ProblemPhyXMini0508.modeNumber_eq_four}
  \uses{def:physics:phyx-mini-0508:phyxminiproblems-problemphyxmini0508-rigidboxelectronsetup, def:physics:phyx-mini-0508:phyxminiproblems-problemphyxmini0508-matchessuppliedwavefunctionfigure, def:physics:phyx-mini-0508:phyxminiproblems-problemphyxmini0508-representsrigidboxstationarymode}
  Four alternating lobes in the supplied plot identify the fourth stationary mode. This is derived from figure evidence and the lobe-count interpretation, not assumed as a raw numerical mode readout
\end{lemma}
\begin{proof}
  The conclusion follows from the governing relations and figure readouts named in the statement of mode Number eq four.
\end{proof}
\begin{lemma}[box Length within one hundredth nanometer]
  \label{lem:physics:phyx-mini-0508:phyxminiproblems-problemphyxmini0508-boxlength-within-one-hundredth-nanometer}
  \lean{PhyXMiniProblems.ProblemPhyXMini0508.boxLength_within_one_hundredth_nanometer}
  \uses{def:physics:phyx-mini-0508:phyxminiproblems-problemphyxmini0508-lengthinnanometers, def:physics:phyx-mini-0508:phyxminiproblems-problemphyxmini0508-rigidboxelectronsetup, def:physics:phyx-mini-0508:phyxminiproblems-problemphyxmini0508-matchesrigidboxelectronscenario, def:physics:phyx-mini-0508:phyxminiproblems-problemphyxmini0508-matchesproblemenergyreadout, def:physics:phyx-mini-0508:phyxminiproblems-problemphyxmini0508-useselectronmassreferencedata, def:physics:phyx-mini-0508:phyxminiproblems-problemphyxmini0508-matchessuppliedwavefunctionfigure, def:physics:phyx-mini-0508:phyxminiproblems-problemphyxmini0508-representsrigidboxstationarymode, def:physics:phyx-mini-0508:phyxminiproblems-problemphyxmini0508-satisfiesinfiniterigidboxenergylaw, def:physics:phyx-mini-0508:phyxminiproblems-problemphyxmini0508-hasphysicalrigidboxparameters}
  With n = 4, E = 6.0 eV, the electron mass, and Physlib's ℏ, the infinite-well law gives approximately 1.00137 nm. The following tolerance is substantially tighter than the tenth-nanometre precision of the choices
\end{lemma}
\begin{proof}
  The conclusion follows from the governing relations and figure readouts named in the statement of box Length within one hundredth nanometer.
\end{proof}
% --- Archon declaration topology end ---
% --- Archon physics formalization source end ---
